\section{서론}

\subsection{연구 필요성}
페달링에서 다리가 크랭크를 돌리는 힘은 크랭크 각도마다 크게 다르다. 다리와 크랭크가 거의 일직선이 되는 두 자세(데드스팟) 근처에서는 힘을 주어도 크랭크가 거의 돌지 않지만, 원형 체인링은 크랭크를 일정한 빠르기로 돌린다.

비원형 체인링은 크랭크 각도마다 체인을 당기는 팔의 길이를 바꾸어, 자전거 속도가 일정해도 팔이 긴 각도에서는 크랭크를 느리게, 짧은 각도에서는 빠르게 돌린다. 근육은 느리게 줄어들수록 큰 힘을 내므로, 힘이 잘 전달되는 구간에서 크랭크를 늦추면 같은 출력을 더 적은 부담으로 낼 수 있다. 이 원리는 근골격 모형 연구에서 제시되었으나\cite{rankin2008} 측정에서는 최대 출력이 달라지지 않아\cite{leong2017}, 효과의 원인과 크기를 식으로 따라갈 수 있는 모델이 필요하다.

\subsection{연구 목적}
\begin{enumerate}
\item 체인링 모양, 페달 토크, 뒷바퀴 저항으로 크랭크 운동을 기술하는 운동방정식을 정의와 법칙에서 유도한다.
\item 같은 속도와 같은 출력으로 달릴 때, 한 바퀴 중 고관절 부담이 가장 큰 순간을 가장 작게 만드는 체인링 모양을 찾는다. 이때 모양을 미리 한 가지 꼴로 정하지 않고 임의의 모양으로 넓혀 푼다.
\item 모델 안에서 힘이 다리와 구동계에 어떻게 나뉘어 전달되는지 그림으로 나타내고, 선행연구의 힘 분석과 비교한다.
\end{enumerate}
